A fundamental goal of robot learning is to build systems that can solve various manipulation tasks in unconstrained 3D settings. A popular class of learning methods directly processes image(s) viewed from single or multiple cameras. These view-based methods have achieved impressive success on a variety of pick-and-place and object rearrangement tasks~\cite{goyal2022ifor,jang2021bcz,shridhar2021cliport,transporter2021}. However, their success on tasks that require 3D reasoning has been limited. As shown by James et al.~\cite{c2farm} and Shridhar et al.~\cite{peract2022arxiv}, view-based methods struggle at 3D manipulation tasks on RLBench~\cite{james2020rlbench} with less than 2\% success.


\begin{wrapfigure}[15]{r}{0.42\textwidth}
\centering
  \vspace{-10pt}
   \includegraphics[width=0.4\textwidth]{figures/time_vs_perf.pdf}
  \caption{RVT scales and performs better than PerAct on RLBench, achieving on-par performance in 36X less time (same hardware), and 1.26X peak performance.}
  \label{fig:rvt_vs_peract}
\end{wrapfigure}


To address this, methods have been proposed that reason with explicit 3D representations of the scene. C2F-ARM~\cite{c2farm} represents the scene with multi-resolution voxels and achieves strong performance on difficult RLBench tasks. PerAct~\cite{peract2022arxiv} improves upon C2F-ARM in behavior cloning by using perceiver transformer~\cite{jaegle2021perceiver} to process voxels. However, creating and reasoning over voxels comes at a higher computing cost compared to reasoning over images, since the number of voxels scales cubicly with the resolution as opposed to squarely for image pixels. This makes voxel-based methods less scalable compared to their view-based counterparts. In fact, training PerAct on 18 RLBench tasks takes 16 days using 8 V100 GPUs (3072 GPU hours). This hinders fast development and prototyping. Moreover, such computing requirements become even more prohibitive when scaling to larger datasets with more tasks and diversity.

Hence, a key question is -- can we build a manipulation network that not only performs well but also inherits the scalability of view-based methods? To this end, we propose \method (\textbf{R}obotic \textbf{V}iew \textbf{T}ransformer) that significantly outperforms the SOTA voxel-based method both in terms of success rate and training time, as shown in Fig.~\ref{fig:rvt_vs_peract}. With the same hardware, \method achieves the peak performance of PerAct in \textbf{36X less time}, decreasing the training time from 14 days to just 10 hours. Apart from being much faster to train, \method also achieves a $\textbf{26\%}$ \textbf{higher success rate} than PerAct, averaged over 18 tasks (249 task variations) on RLBench. \method outperforms PerAct on $\textbf{88.9\%}$ \textbf{of tasks} on RLBench while achieving \textbf{2.3X the inference speed} (11.6 vs 4.9 fps). Further, we find that \method also works well in the real world, where with only 51 demonstrations, a single \method model can learn to perform a variety of manipulation tasks (5 tasks, 13 variations) like opening a drawer, placing objects on a shelf, pressing hand sanitizer, and stacking objects (see Fig.~\ref{fig:real}).

At its core, \method is a view-based method that leverages the transformer architecture. It jointly attends over multiple views of the scene and aggregates information across the views. It then produces view-wise heatmaps and features that are used to predict robot end-effector pose. We extensively explore the design of the multi-view architecture and report several useful findings. For example, we observe a better performance when enforcing the transformer to first attend over patches within the same image before concatenating the patches for joint attention. 

Another key innovation is that, unlike prior view-based methods, we decouple the camera images from the images fed to the transformer, by re-rendering the images from \textit{virtual views}. This allows us to control the rendering process and leads to several benefits. For example, we can re-render from viewpoints that are useful for the task (e.g., directly above the table) while not being restricted by real-world physical constraints. Also, since the multi-view input to \method is obtained via re-rendering, we can use \method even with a single sensor camera -- as done in our real-world experiments.

To summarize, our contributions are threefold: first, we propose \method, a multi-view transformer for 3D object manipulation that is accurate and scalable; second, we investigate various design choices for the multi-view transformer that lead to better object manipulation performance; and finally, we present an empirical study for multi-task object manipulation in simulation and the real world.